\documentclass[10pt,a4paper]{amsart}
\usepackage[margin=19mm]{geometry}
\usepackage{amsmath,amssymb,mathtools}
\usepackage[scr=rsfs,cal=euler]{mathalfa}
\usepackage{xfrac}
\usepackage[unicode,hidelinks,bookmarks=false]{hyperref}
\newcommand{\ie}{i.e.\ }
\newcommand{\cf}{cf.\ }
\newcommand{\resp}{respectively\ }
\newcommand{\Subsection}{\subsection}
\providecommand{\nobreakdash}{\nobreak\hbox{-}}
\setlength{\emergencystretch}{2em}
\multlinegap0pt
\usepackage{latexsym,color}
\usepackage{mathrsfs}
\newtheorem{thm}{Theorem}[section]
\newtheorem{cor}[thm]{Corollary}
\newtheorem{prop}[thm]{Proposition}
\newtheorem{lem}[thm]{Lemma}
\newtheorem{rem}[thm]{Remark}
\newtheorem{defn}[thm]{Definition}
\def\GA{\mathcal A}
\def\CC{\mathbb C}
\def\GC{\mathcal C}
\def\DD{\mathbb D}
\def\GD{\mathcal D}
\def\GF{\mathcal F}
\def\G{\mathcal G}
\def\GH{\mathcal H}
\def\GJ{\mathcal J}
\def\GL{\mathcal L}
\def\GM{\mathcal M}
\def\GN{\mathcal N}
\def\NN{\mathbb N}
\def\GR{\mathcal R}
\def\RR{\mathbb R}
\def\TT{\mathbb T}
\def\GX{\mathcal X}
\def\ZZ{\mathbb Z}
\def\beginpf{\begin{proof}}
\def\endpf{\end{proof}}
\def\beq{\begin{equation}}
\def\eeq{\end{equation}}
\def\ol{\overline}
\def\til{\tilde}
\def\q{\quad}
\def\rank{\mathop{\rm rank}\nolimits}
\def\sigab{\Sigma_{a,b}}
\def\wt{\widetilde}
\def\wX{\widetilde X}
\def\wY{\widetilde Y}
\def\ker{\mathop{\rm ker}\nolimits}
\def\ran{\mathop{\rm ran}\nolimits}
\newcommand{\simast}{%
    \ensuremath{%
       \stackrel{\mathsf{\ast}}{\sim}}}
\def\spam{\mathop{\rm span}\nolimits}   % \span is already in use
\def\ind{\mathop{\rm ind}\nolimits}
\def\CP{{\rm CP}}
\def\GCD{{\rm GCD}}
\def\essinf{\mathop{\rm ess \, inf}}
\def\clos{\mathop{\rm clos}\nolimits}
\begin{document}
\title[Equivalence after extension for general Toeplitz operators]{Equivalence after extension for general Toeplitz operators}
\author{M. Cristina Câmara; Jonathan R. Partington}
\date{}
\maketitle
\noindent\emph{Source and licence.} Equivalence after extension for general Toeplitz operators. Version arXiv:2608.08196v1 (CC BY 4.0). This material is distributed under the Creative Commons Attribution 4.0 International licence, http://creativecommons.org/licenses/by/4.0/, as recorded in the arXiv OAI licence metadata for this exact version and shown on the abstract page https://arxiv.org/abs/2608.08196v1. Full rights notice: licenses/arxiv-2608.08196-CC-BY-4.0.txt.\par\smallskip
\noindent\textit{Editorial scope.} Counted unit: the original \texttt{thm} environment \texttt{thm:2.1aug3} at source lines 255--259 together with its complete proof at lines 260--278; the delivered document keeps source lines 132--281 so that every referenced label and every citation resolves inside it. Native editable source is the paper's own concatenated TeX; the AMS wrapper is editorial formatting only. No publisher class, upstream script or unrelated artwork is redistributed.
\beq\label{eq:4}
G= G_- z^k G_+
\eeq
with $k \in \ZZ$, $G^{\pm 1}_- \in \ol{H^\infty} \cap C^\mu(\TT)$, and
$G^{\pm 1}_+ \in H^\infty \cap C^\mu(\TT)$.
Here $k$ is the winding number of $G$ around the origin in the complex plane.\\

Now denote by $P^+$ the orthogonal projection from $L^2:=L^2(\TT)$ onto $H^2_+:= H^2(\DD)$ and define the {\em Toeplitz operator with symbol $G$\/} by
\beq \label{eq:5}
T_G = P^+ G P^+_{| H^2_+}.
\eeq
The factorization \eqref{eq:4}
induces the operator factorization
\[
T_G = T_{G_-} T_{z^k} T_{G_+},
\] 
where $T_{G_-}$ and $T_{G_+}$ are invertible, so we have that
\[
T_G \sim T_{z^k},
\]
thus reducing the study of many properties of $T_G$ to those of a very simple Toeplitz operator $T_{z^k}$.\\

The concept of equivalent operators was generalised by Bart and Tseka\-novski\u\i~\cite{BTsk} in 1992. We say that two operators $T: X \to \wX$ and $S:Y \to \wY$ are
{\em equivalent after extension\/} (EAE) if and only if there exist Banach spaces $X_0$, $Y_0$
and invertible operators $E,F$ such that
\beq\label{eq:8}
\begin{pmatrix}
T & 0 \\ 0 & I_{X_0}
\end{pmatrix} = 
E
\begin{pmatrix}
S & 0 \\ 0 & I_{Y_0}
\end{pmatrix}
F.
\eeq
The spaces $X_0$ and $Y_0$ are called {\em extension spaces}. Then we write $T \simast S$.
That is,
\[ T \simast S \iff \begin{pmatrix}
T & 0 \\ 0 & I_{X_0}
\end{pmatrix} \sim
\begin{pmatrix}
S & 0 \\ 0 & I_{Y_0}
\end{pmatrix}.
\]
Some shared properties of operators that are EAE are the following:
\begin{itemize}
\item $\ker T \cong \ker S$;
\item $\ran T$ is closed if and only if $\ran S$ is closed, and in that case\\
$\wX/\ran T \cong \wY/\ran S$;
\item if one of the operators is left/right invertible, then so is the other;
\item if $T$ is Fredholm  (that is, if $T$  has closed range and
the spaces   
$\ker T$ and $\wX/\ran T$ are finite-dimensional), then  $S$ is also Fredholm, with $\dim \ker T=\dim \ker S$
and $\dim \wX/\ran T = \dim  \wY/\ran S$.
\end{itemize}
Some other properties are not necessarily shared.
For instance, one operator may be the zero operator, when the other is not; 
moreover, they do not necessarily share the same spectrum.

However, one may reformulate the questions of the zero operator or 
the spectrum in terms of kernels and invertibility, in such a way that one still obtains
answers for one operator from the study of another operator that is EAE. We give examples later.

Useful as the notion of EAE may be, behind it lie several difficult questions.\\

\noindent {\bf Question 1}. How do we show that two operators are EAE?\\

In general we have to construct the operators $E$ and $F$ in \eqref{eq:8} and choose the spaces $X_0$ and $Y_0$. Moreover,
given an operator $T$ whose properties one wants to study, many operators can be chosen that
are EAE to $T$. The choice is useful if studying the corresponding properties
of $S$ is simpler that the original task for $T$. So there are two more natural questions:\\

\noindent {\bf Question 2.} How do we find a {\em useful\/} EAE relation?\\

\noindent {\bf Question 3}. Which extension spaces should we choose?\\

In what follows we take advantage of the strong connection that must exist between the solutions of the two equations
\[
Tf=g, \qquad S\phi=\psi,
\]
and, in particular,
\[
Tf=0, \qquad S\phi=0,
\]
to suggest answers to the questions above in various examples.

\section{Toeplitz operators on $H^2_+$}

Recall that the Toeplitz operator $T_G: H^2_+ \to H^2_+$ with symbol $G \in L^\infty$
was defined  in \eqref{eq:5}.
Its kernel consists of the functions $\phi_+ \in H^2_+$ such that 
$P^+ G \phi_+=0$; that is, $G\phi_+ \in H^2_-:= L^2 \ominus H^2_+$. Thus,
for any $\phi_+ \in H^2_+$ we have
\begin{align}\label{eq:18}
T_G \phi_+ = 0 & \iff G\phi_+\in H^2_-\nonumber
\\
& \iff G\phi_+ + \phi_- = 0 \quad \hbox{for some } \phi_- \in H^2_- \nonumber \\
& \iff (GP^+ + P^-)\phi=0 \quad \hbox{with } \phi \in L^2 \hbox{ and } P^\pm\phi=\phi_\pm.
\end{align}
The operator $GP^+ + P^-$ is called a {\em paired operator}. These are operators
of the form 
\[
S_{A,B} = AP^+ + B P^-,
\]
where $A$ and $B$ are bounded operators on $L^2$. 
For a recent study of their properties we refer to \cite{CGP24,CP24}.
In this case
we take $A$ to be multiplication by $G$ and $B=I$, the identity. From
\eqref{eq:18} we see that
\beq\label{eq:19}
\phi_+ \in \ker T_G \iff \phi \in \ker (GP^++P^-), \hbox{ and } P^+\phi=\phi_+ 
\eeq
and
\[
P^+: \ker (GP^++P^-) \to \ker T_G
\]
is an isomorphism with inverse defined by
\[
\phi_+ \mapsto \phi_+ - G \phi_+ \quad \hbox{for } \phi_+ \in \ker T_G.
\]
From \eqref{eq:19} and $\ker T_G=P^+ \ker (GP^++P^-)$ we get a natural candidate for
an operator that is EAE to $T_G$, as well as a 
hint of which the unknown spaces should be. Indeed, we have the following:

\begin{thm}\label{thm:2.1aug3}
\beq\label{eq:22}
GP^++P^- \simast T_G
\eeq
\end{thm}

\begin{proof}
This follows since
\[
\begin{pmatrix}
T_G & 0 \\ 0 & I_{H^2_-}
\end{pmatrix}
=
\begin{pmatrix}
P^+ & -P^- \\ P^- & P^+ 
\end{pmatrix}
\begin{pmatrix}
GP^++P^- & 0 \\ 0 & I_{\{0\}}
\end{pmatrix}
\begin{pmatrix}
P^+ -P^-GP^+ & P^- \\ 0 & I_{\{0\}}
\end{pmatrix},
\]
where the factors on the left and right of the right-hand side are invertible operators.
\end{proof}

Regarding Question 2, the relation \eqref{eq:22} is useful because there is a well-developed
factorisation theory to study properties such as invertibility and Fredholmness of paired operators.


\begin{thebibliography}{99}
\bibitem[BTsk]{BTsk} Bart, H. and Tsekanovski\u\i, V.\`E.: Matricial coupling and equivalence after extension. In: Ando, T., Gohberg. I. (eds.) Operator theory and complex analysis (Sapporo, 1991), pp. 143--160. Springer, Basel (1992).
\bibitem[CGP24]{CGP24} C\^amara, M. Cristina; Guimar\~{a}es, Andr\'e; Partington, Jonathan R. Paired operators and paired kernels. {\em Postpandemic operator theory}, 15--28, Theta Ser. Adv. Math., 27, Theta, Bucharest, 2024.
\bibitem[CP24]{CP24} C\^amara, M. Cristina; Partington, Jonathan R. Paired kernels and their applications. {\em Results Math.} 79 (2024), no. 3, Paper No. 120, 23 pp.
\end{thebibliography}
\end{document}
